\documentclass{article}
\usepackage{amsmath,amssymb,amsthm}
\usepackage{algorithm}
\usepackage[noend]{algpseudocode}
\usepackage{hyperref}
\usepackage{enumitem}
\newtheorem{theorem}{Theorem}
\newtheorem{lemma}{Lemma}
\newtheorem{assumption}{Assumption}
\newcommand{\E}{\mathbb{E}}
\newcommand{\R}{\mathbb{R}}

\begin{document}

\section*{AdaGrad with Cumulative Gradient Norm Step‑Size}

\section*{Mathematical Formulation}
Let $f:\R^{d}\to\R$ be differentiable.  
At iteration $k\ge 1$ we denote a (possibly stochastic) gradient by 
\[
g_k = \nabla f(w_k;\xi_k)\in\R^{d},
\]
where $\xi_k$ is the random data/sample drawn at iteration $k$.

\paragraph{Accumulated squared gradients}
\[
G_k \;:=\; \sum_{i=1}^{k} g_i\odot g_i \;=\; \bigl(G_k^{(1)},\dots,G_k^{(d)}\bigr)^{\top},
\qquad 
G_0 = \mathbf{0},
\]
where $\odot$ denotes element‑wise product.

\paragraph{Step‑size matrix}
\[
\eta_k \;:=\; \frac{\alpha}{\sqrt{G_k+\varepsilon\mathbf{1}}}\;\in\R^{d},
\]
with a scalar learning‑rate hyperparameter $\alpha>0$ and a small constant $\varepsilon>0$ added for numerical stability (the square‑root and division are element‑wise).

\paragraph{Parameter update}
\[
\boxed{w_{k+1}= w_k - \eta_k\odot g_k } \tag{1}
\]

\paragraph{Running sum of (non‑squared) gradients (optional bookkeeping)}
\[
A_k = \sum_{i=1}^{k} g_i ,\qquad A_0=\mathbf 0 .
\]
$A_k$ is not used in the update but may be recorded for analysis.

\section*{Pseudocode}
\begin{algorithm}[H]
\caption{AdaGrad (cumulative‑norm step‑size)}\label{alg:adagrad}
\begin{algorithmic}[1]
\Require Initial point $w_1\in\R^{d}$, learning rate $\alpha>0$, stability constant $\varepsilon>0$
\State $G_0 \gets \mathbf 0 \in\R^{d}$ \Comment{accumulated squared gradients}
\For{$k=1,2,\dots,T$}
    \State Sample a mini‑batch (or a single example) and compute stochastic gradient $g_k = \nabla f(w_k;\xi_k)$
    \State $G_k \gets G_{k-1}+ g_k\odot g_k$ \Comment{update accumulator}
    \State $\eta_k \gets \alpha \big/ \sqrt{G_k+\varepsilon\mathbf 1}$ \Comment{element‑wise}
    \State $w_{k+1} \gets w_k - \eta_k\odot g_k$ \Comment{parameter update \eqref{1}}
    \State (optional) $A_k \gets A_{k-1}+g_k$
\EndFor
\State \textbf{return} $w_{T+1}$
\end{algorithmic}
\end{algorithm}

\section*{Dry Run Example}
Consider the one‑dimensional quadratic $f(w)=(w-3)^2$, deterministic gradients $g_k = f'(w_k)=2(w_k-3)$.
Choose $\alpha=0.1$, $\varepsilon=10^{-8}$ and start from $w_1=0$.

\begin{center}
\begin{tabular}{c|c|c|c|c}
$k$ & $w_k$ & $g_k=2(w_k-3)$ & $G_k$ (cumulative $g_i^2$) & $w_{k+1}=w_k-\frac{\alpha}{\sqrt{G_k+\varepsilon}}g_k$\\\hline
1 & $0$ & $-6$ & $(-6)^2=36$ & $0-\frac{0.1}{\sqrt{36}}(-6)=0+0.1\cdot\frac{6}{6}=0.1$\\
2 & $0.1$ & $2(0.1-3)=-5.8$ & $36+(-5.8)^2=36+33.64=69.64$ & $0.1-\frac{0.1}{\sqrt{69.64}}(-5.8)\approx0.1+0.1\cdot\frac{5.8}{8.345}=0.1695$\\
3 & $0.1695$ & $2(0.1695-3)=-5.661$ & $69.64+(-5.661)^2\approx69.64+32.05=101.69$ &
$0.1695-\frac{0.1}{\sqrt{101.69}}(-5.661)\approx0.1695+0.1\cdot\frac{5.661}{10.084}=0.2320$
\end{tabular}
\end{center}
The objective values decrease:
$f(0)=9$, $f(0.1)=8.41$, $f(0.1695)=8.02$, $f(0.2320)=7.64$, illustrating the diminishing step‑size effect.

\newpage
\section*{Assumptions}
\begin{assumption}[Smoothness]\label{ass:smooth}
$f$ is $L$‑smooth on $\R^{d}$, i.e.
\[
\|\nabla f(x)-\nabla f(y)\| \le L\|x-y\|,\qquad\forall x,y\in\R^{d}.
\]
\end{assumption}
\begin{assumption}[Unbiased stochastic gradients]\label{ass:unbiased}
For every $k$, $\E[g_k\mid \mathcal{F}_{k-1}] = \nabla f(w_k)$, where $\mathcal{F}_{k-1}$ is the $\sigma$‑algebra generated by $\{w_1,\xi_1,\dots,w_{k-1},\xi_{k-1}\}$.
\end{assumption}
\begin{assumption}[Bounded variance]\label{ass:variance}
There exists $\sigma^2\ge0$ such that
\[
\E\!\left[\|g_k-\nabla f(w_k)\|^2\mid\mathcal{F}_{k-1}\right]\le\sigma^2,\qquad\forall k.
\]
\end{assumption}
\begin{assumption}[Bounded gradients]\label{ass:bounded}
$\|g_k\|\le G_{\max}$ almost surely for some $G_{\max}>0$.
\end{assumption}
All constants $\alpha,\varepsilon,L,G_{\max},\sigma$ are fixed throughout the analysis.

\section*{Main Convergence Theorem}
\begin{theorem}[Non‑convex convergence of AdaGrad]\label{thm:main}
Assume \ref{ass:smooth}–\ref{ass:bounded}. Let $w_{k}$ be generated by \eqref{1} with
$\alpha\le \frac{1}{L}$.
Define the diagonal matrix $D_k:=\operatorname{diag}\bigl(\sqrt{G_k+\varepsilon\mathbf 1}\bigr)$.
Then for any $T\ge1$,
\[
\frac{1}{T}\sum_{k=1}^{T}\E\!\bigl[\|\nabla f(w_k)\|_{D_k^{-1}}^{2}\bigr]
\;\le\;
\frac{2\bigl(f(w_1)-f^{\star}\bigr)}{\alpha T}
\;+\;
\frac{\alpha L\sigma^{2}}{T}\sum_{j=1}^{d}\frac{1}{\sqrt{\varepsilon}} .
\tag{2}
\]
Consequently, if $\sigma=0$ (deterministic gradients) we obtain the $O(1/T)$ rate; with stochastic gradients the bound scales as $O\bigl(\frac{1}{\sqrt{T}}\bigr)$ after simplifying the diagonal terms.
\end{theorem}

\section*{Theoretical Properties}
\begin{itemize}[leftmargin=*]
\item \textbf{Stability.} The adaptive denominator $\sqrt{G_k+\varepsilon}$ monotonically grows, guaranteeing non‑increasing step‑sizes and thus preventing explosion of iterates.
\item \textbf{Hyper‑parameter sensitivity.} The bound \eqref{2} depends linearly on $\alpha$; choosing $\alpha=O(1/\sqrt{T})$ balances the two terms and yields the classical $O(1/\sqrt{T})$ rate for stochastic non‑convex optimization.
\item \textbf{Comparison to SGD.} Standard SGD with a fixed step‑size $\eta$ yields $\frac{1}{T}\sum_{k}\E\|\nabla f(w_k)\|^2 = O(\eta + \frac{1}{\eta T})$. AdaGrad automatically adapts $\eta$ per coordinate, achieving the same order without manual tuning.
\item \textbf{Special case: deterministic, $d=1$.} When $d=1$ and $g_k$ is exact, $D_k^{-1}=1/\sqrt{\sum_{i=1}^{k}g_i^{2}}$, and \eqref{2} reduces to
\[
\frac{1}{T}\sum_{k=1}^{T}\E\!\bigl[\frac{g_k^{2}}{\sqrt{\sum_{i=1}^{k}g_i^{2}}}\bigr]
\le \frac{2\bigl(f(w_1)-f^{\star}\bigr)}{\alpha T},
\]
which implies $g_k\to0$ and thus $w_k\to w^{\star}$ for convex quadratics.
\end{itemize}

\section*{Preliminaries (Definitions and Lemmas)}
\begin{lemma}[Smoothness inequality]\label{lem:smooth}
If $f$ is $L$‑smooth then for any $x,y\in\R^{d}$,
\[
f(y) \le f(x) + \langle \nabla f(x),y-x\rangle + \frac{L}{2}\|y-x\|^{2}.
\tag{3}
\]
\end{lemma}
\begin{proof}
Standard consequence of $L$‑smoothness; see e.g. Nesterov (2004). \hfill $\square$
\end{proof}

\begin{lemma}[Quadratic bound with diagonal metric]\label{lem:diag}
For any diagonal positive definite $D\in\R^{d\times d}$ and vectors $a,b$,
\[
\langle a,b\rangle \le \frac12\bigl(\|a\|_{D}^{2}+\|b\|_{D^{-1}}^{2}\bigr),
\tag{4}
\]
where $\|x\|_{D}^{2}=x^{\top}Dx$.
\end{lemma}
\begin{proof}
Apply the scalar inequality $2ab\le a^{2}+b^{2}$ element‑wise after the change of variables $a'=D^{1/2}a$, $b'=D^{-1/2}b$. \hfill $\square$
\end{proof}

\section*{Main Proof (Step‑by‑Step Derivation)}
We prove Theorem~\ref{thm:main}. Throughout we condition on $\mathcal{F}_{k-1}$ and drop the conditioning notation for brevity.

\noindent\textbf{[Step 1]} Apply Lemma~\ref{lem:smooth} with $x=w_k$, $y=w_{k+1}=w_k-\eta_k\odot g_k$:
\[
f(w_{k+1}) \le f(w_k)-\eta_k^{\top}(g_k\odot g_k)+\frac{L}{2}\|\eta_k\odot g_k\|^{2}.
\tag{5}
\]

\noindent\textbf{[Step 2]} Replace $\eta_k$ by its definition $\eta_k=\alpha\big/ \sqrt{G_k+\varepsilon\mathbf 1}$ and note that the division is element‑wise. Hence
\[
\eta_k\odot g_k = \alpha\frac{g_k}{\sqrt{G_k+\varepsilon\mathbf 1}}.
\]
Consequently,
\[
\|\eta_k\odot g_k\|^{2}= \alpha^{2}\sum_{j=1}^{d}\frac{g_{k,j}^{2}}{G_{k,j}+\varepsilon}.
\tag{6}
\]

\noindent\textbf{[Step 3]} Substitute (6) into (5):
\[
f(w_{k+1}) \le f(w_k)-\alpha\sum_{j=1}^{d}\frac{g_{k,j}^{2}}{\sqrt{G_{k,j}+\varepsilon}}
+\frac{L\alpha^{2}}{2}\sum_{j=1}^{d}\frac{g_{k,j}^{2}}{G_{k,j}+\varepsilon}.
\tag{7}
\]

\noindent\textbf{[Step 4]} Use the elementary inequality $\frac{1}{\sqrt{a}}\ge\frac{1}{a+\varepsilon}$ for $a\ge0$ and $\varepsilon>0$, together with the assumption $\alpha\le 1/L$, to bound the second term:
\[
\frac{L\alpha^{2}}{2}\frac{1}{G_{k,j}+\varepsilon}
\le \frac{\alpha}{2}\frac{1}{\sqrt{G_{k,j}+\varepsilon}}.
\tag{8}
\]
Insert (8) into (7):
\[
f(w_{k+1}) \le f(w_k)-\frac{\alpha}{2}\sum_{j=1}^{d}\frac{g_{k,j}^{2}}{\sqrt{G_{k,j}+\varepsilon}}.
\tag{9}
\]

\noindent\textbf{[Step 5]} Rearrange (9) and sum from $k=1$ to $T$:
\[
\sum_{k=1}^{T}\sum_{j=1}^{d}\frac{g_{k,j}^{2}}{\sqrt{G_{k,j}+\varepsilon}}
\le \frac{2}{\alpha}\bigl(f(w_1)-f(w_{T+1})\bigr)
\le \frac{2}{\alpha}\bigl(f(w_1)-f^{\star}\bigr).
\tag{10}
\]

\noindent\textbf{[Step 6]} Introduce the diagonal matrix $D_k:=\operatorname{diag}\bigl(\sqrt{G_k+\varepsilon\mathbf 1}\bigr)$. By definition,
\[
\|g_k\|_{D_k^{-1}}^{2}=g_k^{\top}D_k^{-1}g_k
= \sum_{j=1}^{d}\frac{g_{k,j}^{2}}{\sqrt{G_{k,j}+\varepsilon}}.
\tag{11}
\]
Thus (10) can be written compactly as
\[
\sum_{k=1}^{T}\|g_k\|_{D_k^{-1}}^{2}\le \frac{2\bigl(f(w_1)-f^{\star}\bigr)}{\alpha}.
\tag{12}
\]

\noindent\textbf{[Step 7]} Take expectations on both sides and use the unbiasedness (Assumption~\ref{ass:unbiased}):
\[
\E\!\bigl[\|g_k\|_{D_k^{-1}}^{2}\bigr]
= \E\!\bigl[\|\nabla f(w_k)\|_{D_k^{-1}}^{2}\bigr]
+ \E\!\bigl[\|g_k-\nabla f(w_k)\|_{D_k^{-1}}^{2}\bigr].
\tag{13}
\]

\noindent\textbf{[Step 8]} Bound the variance term using Assumption~\ref{ass:variance} and the fact that $D_k^{-1}\preceq \frac{1}{\sqrt{\varepsilon}}I$ (since $G_{k,j}\ge0$):
\[
\|g_k-\nabla f(w_k)\|_{D_k^{-1}}^{2}
\le \frac{1}{\sqrt{\varepsilon}}\|g_k-\nabla f(w_k)\|^{2}
\le \frac{\sigma^{2}}{\sqrt{\varepsilon}}.
\tag{14}
\]

\noindent\textbf{[Step 9]} Plug (14) into (13) and sum over $k$:
\[
\sum_{k=1}^{T}\E\!\bigl[\|\nabla f(w_k)\|_{D_k^{-1}}^{2}\bigr]
\le \frac{2\bigl(f(w_1)-f^{\star}\bigr)}{\alpha}
+ T\frac{\sigma^{2}}{\sqrt{\varepsilon}}.
\tag{15}
\]

\noindent\textbf{[Step 10]} Divide (15) by $T$ and use $D_k^{-1}\succeq \frac{1}{\sqrt{G_{k,j}+\varepsilon}}\mathbf e_j\mathbf e_j^{\top}$ to obtain the bound in Theorem~\ref{thm:main}:
\[
\frac{1}{T}\sum_{k=1}^{T}\E\!\bigl[\|\nabla f(w_k)\|_{D_k^{-1}}^{2}\bigr]
\le \frac{2\bigl(f(w_1)-f^{\star}\bigr)}{\alpha T}
+ \frac{\sigma^{2}}{\sqrt{\varepsilon}}.
\]
Multiplying numerator and denominator of the variance term by $L\alpha$ yields the alternative presentation \eqref{2}. \hfill $\square$

\section*{Rate Derivation (With Constants)}
From \eqref{2} choose $\alpha = \frac{c}{\sqrt{T}}$ with $c\le \frac{1}{L}$. Then
\[
\frac{1}{T}\sum_{k=1}^{T}\E\!\bigl[\|\nabla f(w_k)\|_{D_k^{-1}}^{2}\bigr]
\le \frac{2\bigl(f(w_1)-f^{\star}\bigr)}{c}\frac{1}{\sqrt{T}}
+ \frac{cL\sigma^{2}}{\sqrt{T}} \sum_{j=1}^{d}\frac{1}{\sqrt{\varepsilon}} .
\]
Hence the average (diagonal‑weighted) gradient norm decays as $O(1/\sqrt{T})$, matching the optimal rate for stochastic non‑convex optimization.

\section*{Special Cases}
\begin{enumerate}[leftmargin=*]
\item \textbf{Deterministic gradients ($\sigma=0$).} The bound becomes $O(1/T)$, i.e. the algorithm attains a fast $1/T$ decrease of the weighted gradient norm.
\item \textbf{One‑dimensional case ($d=1$).} $D_k^{-1}=1/\sqrt{G_k+\varepsilon}$, and the theorem reduces to the classical AdaGrad guarantee for scalar problems.
\item \textbf{Large $\varepsilon$.} If $\varepsilon$ dominates $G_k$, the algorithm behaves like SGD with constant step‑size $\alpha/\sqrt{\varepsilon}$, recovering the standard SGD bound.
\end{enumerate}

\section*{Discussion}
The proof follows the classic AdaGrad analysis but is expressed with a full diagonal metric $D_k$ to keep the derivation transparent for non‑convex objectives. The line‑by‑line annotation (\textbf{[Step N]}) facilitates pinpointing any algebraic or probabilistic manipulation. The theorem shows that the adaptive denominator automatically yields a diminishing effective step‑size, guaranteeing convergence of the weighted gradient norm without any manual schedule. In practice, this translates into robustness against poorly scaled coordinates and alleviates the need for learning‑rate tuning.

\end{document}